\subsection{两个自同态的张量积的特征多项式}

命题 11. --- 设 E（相应地 E'）是交换域 K 上的有限维向量空间，并设 u（相应地 u'）是 E（相应地 E'）的自同态。设

\[
\chi_ {u} (\mathbf {X}) = \prod_ {i} \left(\mathbf {X} - \alpha_ {i}\right), \quad \chi_ {u ^ {\prime}} (\mathbf {X}) = \prod_ {j} \left(\mathbf {X} - \beta_ {j}\right)
\]

是 u 与 \(u'\) 的特征多项式在 K 的某个适当扩域中的线性因式分解。则自同态 \(u \otimes u'\)（向量空间 \(E \otimes_{K} E'\) 的）的特征多项式由以下公式给出

\[
\chi_ {u \otimes u ^ {\prime}} (\mathbf {X}) = \prod_ {i, j} \left(\mathbf {X} - \alpha_ {i} \beta_ {j}\right).
\]

如同 VII，第 36 页，命题 10 的证明中那样论证，我们看到只需证明以下引理：

引理2. --- 设 B 和 C 分别是 m 阶与 n 阶的下三角矩阵，其对角线分别为 \((\beta_{i})_{1\leq i\leq m}\) 与 \((\gamma_{j})_{1\leq j\leq n}\)。把有序集 \(\{1,2,\ldots,m\}\) 与 \(\{1,2,\ldots,n\}\) 的字典序积等同于区间 \(\{1,2,\ldots,mn\}\)。则张量积矩阵（II，第 357 页）\(B\otimes C\) 是下三角矩阵，其对角线为 \((\beta_{i}\gamma_{j})\)。

这直接由两个矩阵的张量积的定义（同上）及字典序乘积的定义（集合论，III，第157页）得出。

